\begin{table}[htbp]
\centering
\caption{Table 2---Effects of Dust Bowl Erosion on Land Values (Simplified TWFE)}
\label{tab:table2}
\begin{tabular}{lcccc}
\toprule
 & Weighted & Unweighted & State$\times$Year FE & High only \\
 & (1) & (2) & (3) & (4) \\
\midrule
\multicolumn{5}{l}{\textit{Dep.\ var.: ln(value farmland + buildings / acre)}} \\
\addlinespace
High erosion $\times$ post-1930 & -0.161*** & -0.205*** & -0.074 & -0.262*** \\
 & (0.058) & (0.054) & (0.108) & (0.053) \\
\addlinespace
Medium erosion $\times$ post-1930 & 0.264*** & 0.223*** & 0.094 & \\
 & (0.052) & (0.046) & (0.080) & \\
\addlinespace
County FE & Y & Y & & Y \\
Year FE & Y & Y & & Y \\
State $\times$ Year FE & & & Y & \\
Weighted (farmland) & Y & & Y & Y \\
\addlinespace
Observations & 11,775 & 11,775 & 11,775 & 11,775 \\
Counties & 785 & 785 & 785 & 785 \\
\bottomrule
\multicolumn{5}{p{0.95\linewidth}}{\footnotesize
 \textit{Notes:} Simplified two-way FE regressions for weight decomposition.
 Full paper uses state$\times$year FE and 1930 baseline controls.
 High/Medium erosion = share of county land in high/medium erosion areas
 (from 1930s Soil Conservation Service maps).
 Robust standard errors clustered by county in parentheses.
 *** Significant at 1\%, ** 5\%, * 10\%.}
\end{tabular}
\end{table}
